\documentclass[12pt,a4paper]{article}
\usepackage{goyabase}

\usepackage{longtable}
\usepackage{calc}

% Pandoc compatibility
\providecommand{\tightlist}{%
  \setlength{\itemsep}{0pt}\setlength{\parskip}{0pt}}
\newcommand{\passthrough}[1]{#1}

\title{\textbf{MongoDB en GloboTour: El Catálogo Flexible}}
\author{Departamento de Informática}
\date{\today}

\usepackage[spanish, es-noshorthands, es-tabla]{babel}
\usepackage[autostyle]{csquotes}
\begin{document}

\maketitle

\section{MongoDB en GloboTour: El Catálogo Flexible}

En \textbf{GloboTour}, el catálogo de experiencias crece cada día.
Tenemos desde cursos de cocina hasta alquiler de yates. Usamos MongoDB
porque nos permite guardar cada experiencia con sus propios campos sin
complicaciones.

\begin{center}\rule{0.5\linewidth}{0.5pt}\end{center}

\subsection{Conceptos Clave (Teoría)}

\begin{longtable}[]{@{}
  >{\raggedright\arraybackslash}p{(\columnwidth - 4\tabcolsep) * \real{0.3333}}
  >{\raggedright\arraybackslash}p{(\columnwidth - 4\tabcolsep) * \real{0.3333}}
  >{\raggedright\arraybackslash}p{(\columnwidth - 4\tabcolsep) * \real{0.3333}}@{}}
\toprule
Concepto SQL & Concepto MongoDB & Descripción \\
\midrule
\endhead
\textbf{Tabla} & \textbf{Colección} & Agrupación de documentos (no requiere esquema fijo). \\
\textbf{Fila} & \textbf{Documento} & Registro individual en formato \textbf{BSON} (JSON binario). \\
\textbf{Columna} & \textbf{Campo} & Par Clave-Valor dentro de un documento. \\
\bottomrule
\end{longtable}

\subsubsection{El dilema del Modelado: ¿Embeber o Referenciar?}

A diferencia de SQL (donde normalizamos), en MongoDB tenemos dos opciones:
\begin{enumerate}
    \item \textbf{Embeber (Denormalizar)}: Metemos los datos dentro del documento (ej: los detalles técnicos dentro de la experiencia). Es más rápido para leer.
    \item \textbf{Referenciar (Normalizar)}: Guardamos un ID de otro documento. Se usa si los datos crecen mucho (ej: miles de comentarios en un post).
\end{enumerate}

\textbf{¡Ojo con los Tipos de Datos!}: En MongoDB, \texttt{String} y \texttt{Number} \textbf{no son lo mismo}. Si guardas el precio como texto, las consultas de \enquote{mayor que} (\texttt{\$gt}) no funcionarán como esperas.

\textbf{Tip de Depuración}: Si tu consulta no devuelve nada:
\begin{enumerate}
    \item Asegúrate de estar en la base de datos correcta: \texttt{globotour}.
    \item Comprueba el tipo de dato real con: \texttt{typeof}.
\end{enumerate}

\begin{center}\rule{0.5\linewidth}{0.5pt}\end{center}

\subsection{Visualización con Mongo Express (GUI)}

\begin{enumerate}
    \item Abre tu navegador en \url{http://localhost:8081}.
    \item \textbf{Login}: Usuario: \texttt{admin} / Password: \texttt{pass}.
    \item \textbf{Navegación}: Haz clic en la base de datos \texttt{globotour} y luego en la colección \texttt{experiencias} para explorar los documentos JSON.
    \item \textbf{Editor Visual}: Puedes editar el JSON directamente desde la web para probar cambios rápidos sin escribir comandos.
\end{enumerate}

\begin{center}\rule{0.5\linewidth}{0.5pt}\end{center}

\subsection{Operaciones Básicas (CRUD)}\label{mql}



\subsubsection{Inserción de Documentos (Create)}
\begin{lstlisting}[language=JavaScript]
// Insertar una nueva experiencia individual
db.experiencias.insertOne({
  nombre: "Cena a ciegas",
  tipo: "Gastronomia",
  precio: 65,
  tags: ["romantico", "sensorial"]
})
\end{lstlisting}

\subsubsection{Consultas de Lectura (Read)}
\begin{lstlisting}[language=JavaScript]
// Sacar todas las experiencias
db.experiencias.find()

// Buscar por coincidencia exacta y operadores
db.experiencias.find({ precio: { $in: [25, 45, 80] } })

// Buscar documentos que TIENEN un campo concreto
db.experiencias.find({ detalles: { $exists: true } })
\end{lstlisting}

\subsubsection{Actualización de Datos (Update)}
\begin{lstlisting}[language=JavaScript]
// Modificar un campo específico ($set)
db.experiencias.updateOne(
  { nombre: "Cena a ciegas" },
  { $set: { precio: 70, dificultad: "Baja" } }
)

// Añadir un elemento a un array ($push)
db.experiencias.updateOne(
  { nombre: "Cena a ciegas" },
  { $push: { tags: "popular" } }
)
\end{lstlisting}

\subsubsection{Borrado de Datos (Delete)}
\begin{lstlisting}[language=JavaScript]
// Borrar un documento concreto
db.experiencias.deleteOne({ nombre: "Cena a ciegas" })
\end{lstlisting}

\begin{center}\rule{0.5\linewidth}{0.5pt}\end{center}

\subsection{Desafíos de Aprendizaje}\label{desafios}

\begin{enumerate}
    \item \textbf{Esquema Flexible}: Inserta una nueva experiencia con campos únicos (por ejemplo, un \texttt{"curso de surf"} con un instructor (\texttt{"Nacho"}) y materiales incluidos (\texttt{material\_incluido: ["tabla","neopreno"]}). Comprueba que MongoDB la acepta sin necesidad de alterar ninguna tabla.
    \item \textbf{Proyecciones}: Modifica una consulta \texttt{find} para que solo devuelva el nombre del producto, ocultando el \texttt{\_id}.
    \item \textbf{Ampliación Avanzada}: Una vez comprendida la base, dirígete al \textbf{02-Laboratorio\_Agregaciones.tex} para dominar el Pipeline de Agregación y el análisis complejo de datos.
\end{enumerate}

\begin{center}\rule{0.5\linewidth}{0.5pt}\end{center}

\subsection{Solucionario de Desafíos}\label{solucionario-mongodb}

\begin{tcolorbox}[title=Solución Desafío 1: Esquema Flexible]
Podemos añadir campos que no existen en otros documentos, como \texttt{instructor} o \texttt{material\_incluido}:
\begin{lstlisting}[language=JavaScript]
db.experiencias.insertOne({
  nombre: "Curso de Surf",
  instructor: "Nacho",
  material_incluido: ["tabla", "neopreno"],
  precio: 40
})
\end{lstlisting}
Si haces un \texttt{db.experiencias.find()}, verás que este documento convive perfectamente con los demás.
\end{tcolorbox}

\begin{tcolorbox}[title=Solución Desafío 2: Proyecciones]
El segundo parámetro de \texttt{find()} permite elegir qué campos mostrar (\texttt{1}) u ocultar (\texttt{0}). El \texttt{\_id} es el único que se muestra siempre a menos que se oculte explícitamente:
\begin{lstlisting}[language=JavaScript]
// Mostrar solo nombre (1) y ocultar id (0)
db.experiencias.find({}, { nombre: 1, _id: 0 })
\end{lstlisting}
\end{tcolorbox}

\end{document}
